\section*{Capítulo \ref{ch:b3:innate-immunity} --- Imunidade inata e inflamação}

\begin{solution}{exo:b3:innate-immunity:1}
\omterm{def:b3:innate-immunity:innate}{Barreiras}: a queratina da pele (mecânica), o muco e os \omterm{def:b3:cytoskeleton-motility:cilia}{cílios} das vias
aéreas (aprisionamento e remoção), o ácido do estômago (mata os
micróbios engolidos), a lisozima e as defensinas das secreções (lisam e
perfuram), o \omterm{def:b3:microbiomes:symbiosis}{microbioma} comensal (ocupa os nichos). Células:
\omterm{def:b3:innate-immunity:innate}{neutrófilos} (matam bactérias por \omterm{def:b3:innate-immunity:killing}{fagocitose}), \omterm{def:b3:innate-immunity:innate}{macrófagos} (comem, fazem
a limpeza, dão o alarme), \omterm{def:b3:innate-immunity:innate}{células dendríticas} (levam o antígeno aos
linfonodos e iniciam as respostas adaptativas), mastócitos (liberam
histamina, inflamam), \omterm{def:b3:innate-immunity:innate}{células natural killer} (matam células infectadas
e transformadas).
\end{solution}

\begin{solution}{exo:b3:innate-immunity:2}
Ser compartilhada por uma classe inteira de micróbios, ser essencial
para eles (de modo que não possa ser perdida para escapar) e estar
ausente no hospedeiro. \omterm{def:b3:bacteriology:envelope}{Lipopolissacarídeo} --- TLR4; flagelina --- TLR5;
RNA de fita dupla --- TLR3 (e RIG-I no citosol); DNA \omterm{def:b3:chromatin-epigenetics:methylation}{CpG} não metilado
--- TLR9; fragmentos de \omterm{def:b3:bacteriology:envelope}{peptidoglicano} --- proteínas NOD;
$\beta$-glicano --- dectina-1.
\end{solution}

\begin{solution}{exo:b3:innate-immunity:3}
Rubor: dilatação arteriolar pela histamina, pelo óxido nítrico e pelas
prostaglandinas. Calor: o mesmo aumento do fluxo sanguíneo. Tumor:
vazamento venular de plasma após a contração endotelial (histamina,
C3a, C5a). Dor: a bradicinina e a prostaglandina E$_{2}$ sensibilizando
os nociceptores, mais a pressão do inchaço.
\end{solution}

\begin{solution}{exo:b3:innate-immunity:4}
Clássica: a C1q ligando-se a anticorpo (ou à \omterm{prop:b3:innate-immunity:systemic}{proteína C reativa}) sobre
uma superfície. Das lectinas: a lectina ligante de manose sobre
açúcares microbianos. Alternativa: a ativação espontânea da C3 e a
deposição de C3b sobre qualquer superfície desprotegida. Desfechos: a
C3b opsoniza para a \omterm{def:b3:innate-immunity:killing}{fagocitose}; a C3a e a C5a inflamam e recrutam; a
C5b inicia o \omterm{def:b3:innate-immunity:complement}{complexo de ataque à membrana}, que lisa.
\end{solution}

\begin{solution}{exo:b3:innate-immunity:5}
$B = [k_{0}/(a-d)](e^{(a-d)t} - 1)$ com $a - d = 0.1$: a \qty{60}{s},
$20(e^{6} - 1) \approx 8000$; a \qty{90}{s}, $20(e^{9} - 1) \approx
1.6\times 10^{5}$. Célula do hospedeiro: $k_{0}/(d - a) = 2/0.1 = 20$.
\end{solution}

\begin{solution}{exo:b3:innate-immunity:6}
Rolamento sobre as \omterm{def:b3:innate-immunity:inflammation}{selectinas} (lectinas que ligam carboidrato);
sinalização por \omterm{def:b3:innate-immunity:inflammation}{quimiocina} que ativa as \omterm{def:b3:innate-immunity:inflammation}{integrinas}; adesão firme das
\omterm{def:b3:innate-immunity:inflammation}{integrinas} às moléculas de adesão endoteliais (ICAM); transmigração
entre as células endoteliais; quimiotaxia ao longo do gradiente. Sem
\omterm{def:b3:innate-immunity:inflammation}{integrinas} $\beta_{2}$ os \omterm{def:b3:innate-immunity:innate}{neutrófilos} rolam mas nunca param nem
atravessam: acumulam-se no sangue (contagens altíssimas), as infecções
se repetem sem formar pus, as feridas cicatrizam mal e o cordão
umbilical cai tarde.
\end{solution}

\begin{solution}{exo:b3:innate-immunity:7}
Doença granulomatosa crônica: a \omterm{def:b3:innate-immunity:killing}{NADPH oxidase} é defeituosa, de modo que
não se faz superóxido, peróxido de hidrogênio nem hipoclorito no
fagossomo. Infecções por organismos que destroem seu próprio peróxido
de hidrogênio (catalase-positivos): \emph{Staphylococcus aureus},
\emph{Aspergillus}, \emph{Burkholderia}, \emph{Serratia},
\emph{Nocardia}. Formam-se granulomas porque os \omterm{def:b3:innate-immunity:innate}{macrófagos} englobam
micróbios que não conseguem matar e os muram, com recrutamento
contínuo, numa lesão crônica.
\end{solution}

\begin{solution}{exo:b3:innate-immunity:8}
Os receptores inibitórios contam o MHC~I; os ativadores contam os
ligantes de estresse; a célula morre se a ativação superar a inibição.
(a) Perder o MHC~I remove a inibição: a \omterm{def:b3:innate-immunity:innate}{célula natural killer} a mata ---
o preço de escapar das células T. (b) Com o MHC~I mantido mas ligantes
de estresse expostos, a ativação ainda pode vencer e a célula é morta;
os \omterm{def:b3:cancer-biology:hallmarks}{tumores}, por isso, muitas vezes descartam ou suprimem os ligantes de
estresse.
\end{solution}

\begin{solution}{exo:b3:innate-immunity:9}
Uma proteína pura não traz nenhum padrão de patógeno, de modo que as
\omterm{def:b3:innate-immunity:innate}{células dendríticas} que a captam não amadurecem, apresentam-na sem
coestimulação, e as células T que a reconhecem são desligadas ou a
ignoram. Um adjuvante --- alume, um lipídio, um ligante de \omterm{def:b3:innate-immunity:prr}{receptor do
tipo Toll} --- dispara os receptores de padrões, amadurece a \omterm{def:b3:innate-immunity:innate}{célula
dendrítica} e fornece o segundo sinal. É o argumento de Janeway: o
sistema adaptativo responde ao antígeno apenas quando o sistema inato
diz que o antígeno veio com um micróbio.
\end{solution}

\begin{solution}{exo:b3:innate-immunity:10}
Sem febre, $48$ duplicações: $10^{4}\times 2^{48} \approx 3\times
10^{18}$ (um absurdo que mostra que quem decide o desfecho é a morte, e
não o crescimento). Com febre, $32$ duplicações: $4\times 10^{13}$ ---
um fator $2^{16} \approx 65\,000$ de bactérias a menos para os
\omterm{def:b3:innate-immunity:innate}{neutrófilos} enfrentarem. Custo: $3\times 0.12\times 8 = \qty{2.9}{MJ}$
por dia, uma refeição farta. Vale a pena quando o patógeno é retardado
e o hospedeiro tem reservas; não vale quando o patógeno cresce igualmente
bem a \qty{40}{\celsius}, ou quando o hospedeiro passa fome, ou quando a
própria febre se aproxima de temperaturas perigosas.
\end{solution}

\begin{solution}{exo:b3:innate-immunity:11}
Ligar o fator H recruta o próprio regulador do hospedeiro: $d$ sobe
acima de $a$ e a alça decai sobre o micróbio como sobre uma célula do
hospedeiro. Clivar a C3b remove a convertase depositada: de novo uma
elevação de $d$. Uma cápsula apresenta uma superfície sobre a qual a
convertase se forma mal, baixando $a$, e esconde dos receptores dos
fagócitos a C3b que é depositada. A mais barata: uma única proteína de
superfície que liga o fator H, que o hospedeiro fornece de graça e que
faz o trabalho.
\end{solution}

\begin{solution}{exo:b3:innate-immunity:12}
Numa farpa, a dilatação, o vazamento e o recrutamento envolvem alguns
mililitros de vasos e entregam defesa onde ela é necessária; os mesmos
mediadores lançados em toda a circulação dilatam todos os vasos e fazem
vazar todos os capilares, de modo que a pressão despenca e a coagulação
é disparada em toda parte --- a fisiologia é idêntica, a escala é
letal. O TNF move a \omterm{def:b3:innate-immunity:inflammation}{inflamação} sinovial da artrite reumatoide, e por
isso bloqueá-lo alivia a doença; mas o TNF é também o que mantém os
\omterm{def:b3:innate-immunity:innate}{macrófagos} organizados nos granulomas que contêm bacilos da tuberculose
dormentes, e bloqueá-lo deixa a tuberculose latente escapar.
\end{solution}

\begin{solution}{pb:b3:innate-immunity:1}
\textbf{1.} Uma hora são $1.5$ duplicações: $10^{4}\times 2^{1.5} \approx
2.8\times 10^{4}$.
\textbf{2.} $2\times 10^{4}$ \omterm{def:b3:innate-immunity:innate}{neutrófilos} matam até $4\times 10^{5}$;
as bactérias crescem apenas de $2.8\times 10^{4}$ a $8\times 10^{4}$:
a morte supera o crescimento em cinco vezes.
\textbf{3.} Hora 2: $8\times 10^{4}$ antes da matança, nenhuma depois. A
população atinge o pico em cerca de $8\times 10^{4}$ perto de
\qty{2}{h} e é zero daí em diante; nenhuma resta a \qty{6}{h}.
\textbf{4.} Chegada a \qty{3}{h}: $10^{4}\times 2^{4.5} = 2.3\times
10^{5}$ nesse momento. Hora 4: $6.4\times 10^{5} - 4\times 10^{5} = 2.4\times
10^{5}$; hora 5: $6.8\times 10^{5} - 4\times 10^{5} = 2.8\times 10^{5}$; hora 6: $7.9\times 10^{5} - 4\times 10^{5} = 3.9\times 10^{5}$ e subindo --- os \omterm{def:b3:innate-immunity:innate}{neutrófilos} já não acompanham, e
forma-se um abscesso. Duas horas de atraso transformam uma cura numa
infecção crônica.
\textbf{5.} $8\times 10^{4}/20 = 4000$ \omterm{def:b3:innate-immunity:innate}{neutrófilos} morrem matando, mais
os $2\times 10^{4}$ que chegaram e morrem de qualquer modo em um dia. O
pus são \omterm{def:b3:innate-immunity:innate}{neutrófilos} mortos, bactérias mortas e tecido liquefeito; os
\omterm{def:b3:innate-immunity:innate}{macrófagos} o removem comendo os cadáveres, ou ele precisa ser drenado.
\textbf{6.} Chegadas de $2000$ por hora, que matam $4\times 10^{4}$:
hora 2, $8\times 10^{4} - 4\times 10^{4} = 4\times 10^{4}$; hora 3,
$1.1\times
10^{5} - 4\times 10^{4} = 7\times 10^{4}$; hora 4, $2\times 10^{5} -
4\times 10^{4} = 1.6\times 10^{5}$ --- crescendo sem limite. O
paciente não consegue conter um inóculo trivial, de modo que os
\omterm{def:b3:bacteriology:antibiotics}{antibióticos} têm de fazer a matança desde o primeiro sinal.
\textbf{7.} $B = 12.5(e^{0.08t} - 1)$: a \qty{30}{s}, $125$; a
\qty{60}{s}, $1500$; a \qty{120}{s}, $1.8\times 10^{5}$.
\textbf{8.} $e^{0.08t} = 1.6\times 10^{5}$: $t = 12/0.08 = \qty{150}{s}$.
\textbf{9.} $k_{0}/(d - a) = 10$.
\textbf{10.} $B(120) = 33(e^{3.6} - 1) \approx 1200$, $150$ vezes menos.
A cápsula compra minutos --- tempo para se multiplicar, e proteção até
que cheguem anticorpos para iniciar a via clássica sobre a própria
cápsula.
\textbf{11.} $10^{4}\times 2\times 10^{6} = 2\times 10^{10}$ C3b: um
milionésimo da C3 do plasma.
\textbf{12.} Sem C3, as três vias param em seu passo comum: nenhuma
\omterm{def:b3:innate-immunity:complement}{opsonização}, nenhuma \omterm{def:b3:innate-immunity:complement}{anafilatoxina}, nenhuma lise --- infecções graves
e recorrentes por bactérias encapsuladas. Sem C9 perde-se apenas o poro
terminal; a \omterm{def:b3:innate-immunity:complement}{opsonização} e a \omterm{def:b3:innate-immunity:inflammation}{inflamação} funcionam, e o fenótipo são
infecções recorrentes por \emph{Neisseria}, o único gênero que os
fagócitos controlam mal e a lise controla bem.
\textbf{13.} $2\times 0.24\times 8 = \qty{3.8}{MJ}$, cerca de
\qty{100}{g} de gordura.
\textbf{14.} Uma duplicação na hora: $2\times 10^{4}$ em vez de
$2.8\times 10^{4}$, um fator $1.4$.
\textbf{15.} Hora 2: $4\times 10^{4}$ antes da matança, nenhuma depois
--- eliminadas a \qty{2}{h}, como antes; a contribuição da febre importa
quando o suprimento de \omterm{def:b3:innate-immunity:innate}{neutrófilos} é marginal, como nas questões 4 e 6.
\textbf{16.} $199\,\text{mg/L}\times 3\,\text{L} = \qty{0.6}{g}$; $0.6/
\num{115000} = 5.2\times 10^{-6}$ mol $= 3.1\times 10^{18}$ moléculas;
ao longo de \qty{172800}{s}, $1.8\times 10^{13}$ por segundo.
\textbf{17.} Eles baixam o ponto de ajuste de volta para perto de
\qty{37}{\celsius}, o paciente se sente melhor e as bactérias dobram um
pouco mais depressa; neste modelo os \omterm{def:b3:innate-immunity:innate}{neutrófilos} ainda vencem, e as
evidências de que os antipiréticos pioram infecções comuns são fracas.
\textbf{18.} Acima de cerca de \qty{41}{\celsius} as proteínas começam a
desnaturar, as enzimas falham e os neurônios disparam errado
(convulsões); \qty{42}{\celsius} é letal. O ponto de ajuste é mantido
abaixo disso por criógenos --- interleucina-10, vasopressina,
glicocorticoides --- e pelo teto da ação das prostaglandinas no
hipotálamo.
\textbf{19.} $10^{9}\times 3\times 10^{6} = 3\times 10^{15}$ moléculas
$= 5\times 10^{-9}$ mol em \qty{5}{L}: \qty{1e-9}{mol/L}.
\textbf{20.} Igual à concentração saturante: todos os \omterm{def:b3:innate-immunity:innate}{macrófagos} do
corpo são ativados ao máximo de uma só vez.
\textbf{21.} $10^{11}\times 1000\times 3600 = 3.6\times 10^{17}$ moléculas
$= 6\times 10^{-7}$ mol; em \qty{15}{L}, \qty{4e-8}{mol/L} --- quatro
mil vezes a concentração ativa.
\textbf{22.} Toda arteríola dilata e toda vênula vaza: o volume de
sangue se acumula nos vasos alargados e escapa para os tecidos, de modo
que a pressão cai; os rins, mal perfundidos, param de filtrar; e a
coagulação disparada em milhares de capilares os obstrui e consome os
fatores de coagulação.
\textbf{23.} As \omterm{def:b3:bacteriology:envelope}{Gram-negativas} mortas liberam de uma vez todo o seu
\omterm{def:b3:bacteriology:envelope}{lipopolissacarídeo}, um bolus para o TLR4, de modo que a onda de
\omterm{def:b3:innate-immunity:inflammation}{citocinas} pode piorar por horas. O anti-TNF falha porque, na hora do
choque, a cascata já passou do TNF --- a interleucina-1, a
interleucina-6 e a coagulação já estão em curso --- e porque o TNF é
necessário para controlar as bactérias: o problema é o momento, e não
o alvo.
\textbf{24.} O \omterm{def:b3:bacteriology:envelope}{lipopolissacarídeo} injetado nada faz ao camundongo sem
TLR4, que sobrevive a doses letais para camundongos normais; uma
infecção viva por \omterm{def:b3:bacteriology:envelope}{Gram-negativas} o mata, porque ele não consegue
detectar as bactérias cedo e montar a \omterm{def:b3:innate-immunity:inflammation}{inflamação} que as contém. O
paradoxo se dissolve: o detector que causa o choque é o detector que
fornece a defesa.
\textbf{25.} Nenhuma bactéria resta a \qty{6}{h} (todas mortas até
\qty{2}{h}); cerca de $1.8\times 10^{5}$ C3b por bactéria a
\qty{120}{s}; dois dias de febre custam \qty{3.8}{MJ}, uns
\qty{100}{g} de gordura.
\end{solution}
